% Idioma.
\usepackage[spanish]{babel}

% Codificación.
\usepackage[T1]{fontenc}
\usepackage[utf8]{inputenc}
\usepackage{bookmark}
\usepackage{float}

% Márgenes.
\usepackage[a4paper,margin=2.5cm]{geometry}

% Matemática.
\usepackage{amsmath}
\usepackage{amssymb}
\usepackage{amsthm}

% Tablas.
\usepackage{booktabs}
\usepackage{longtable}
\usepackage{array}
\usepackage{seqsplit}

\newcolumntype{Y}{>{\raggedright\arraybackslash}X}
\newcolumntype{L}[1]{>{\raggedright\arraybackslash}p{#1}}
\newcolumntype{C}[1]{>{\centering\arraybackslash}p{#1}}
\newcolumntype{R}[1]{>{\raggedleft\arraybackslash}p{#1}}

% Table helpers.
\newcommand{\TablaClave}[1]{\seqsplit{#1}}

\newlength{\TablaColNombre}
\setlength{\TablaColNombre}{0.16\linewidth}

\newlength{\TablaColDescripcion}
\setlength{\TablaColDescripcion}{0.42\linewidth}

\newlength{\TablaColCategoria}
\setlength{\TablaColCategoria}{0.20\linewidth}

\newlength{\TablaColTipo}
\setlength{\TablaColTipo}{0.14\linewidth}

\newcommand{\TablaFila}[4]{#1 & #2 & #3 & #4 \\
\hline}

\newcommand{\TablaFilaFlexible}[1]{#1 \\
\hline}

\newcommand{\TablaCabecera}[4]{%
\toprule
\multicolumn{1}{c}{\textbf{#1}} &
\multicolumn{1}{c}{\textbf{#2}} &
\multicolumn{1}{c}{\textbf{#3}} &
\multicolumn{1}{c}{\textbf{#4}} \\
\midrule
\endfirsthead
\toprule
\multicolumn{1}{c}{\textbf{#1}} &
\multicolumn{1}{c}{\textbf{#2}} &
\multicolumn{1}{c}{\textbf{#3}} &
\multicolumn{1}{c}{\textbf{#4}} \\
\midrule
\endhead
\midrule
\endfoot
\bottomrule
\endlastfoot
}

\newcommand{\TablaCabeceraFlexible}[1]{%
\toprule
#1 \\
\midrule
\endfirsthead
\toprule
#1 \\
\midrule
\endhead
\midrule
\endfoot
\bottomrule
\endlastfoot
}

\NewDocumentEnvironment{TablaDatos}{mmmm}
{
    % More breathing room around the table.
    \setlength{\tabcolsep}{7pt}
    \renewcommand{\arraystretch}{1.8}
    \setlength{\LTleft}{0pt}
    \setlength{\LTright}{0pt}
    \begin{longtable}{
        >{\raggedright\arraybackslash}p{\TablaColNombre}
        >{\raggedright\arraybackslash}p{\TablaColDescripcion}
        >{\raggedright\arraybackslash}p{\TablaColCategoria}
        >{\raggedright\arraybackslash}p{\TablaColTipo}
    }
    \TablaCabecera{#1}{#2}{#3}{#4}
}
{
    \end{longtable}
}

\NewDocumentEnvironment{TablaDatosFlexible}{mm}
{
    % More breathing room around the table.
    \setlength{\tabcolsep}{7pt}
    \renewcommand{\arraystretch}{1.8}
    \setlength{\LTleft}{0pt}
    \setlength{\LTright}{0pt}
    \begin{longtable}{#1}
    \TablaCabeceraFlexible{#2}
}
{
    \end{longtable}
}

% Gráficos.
\usepackage{graphicx}
\usepackage{pgfplots}
\pgfplotsset{compat=1.18}

% Colores.
\usepackage{xcolor}

% Hipervínculos.
\makeatletter
\@ifpackageloaded{hyperref}{%
    \hypersetup{hidelinks}%
}{%
    \usepackage[hidelinks]{hyperref}%
}
\makeatother

% Commands.

% Casi todas las secciones deben tener la subsubsection de análisis descriptivo.
\newcommand{\AnalisisDescriptivo}{\subsubsection{Análisis descriptivo}}

% Tabla de datos.
\newcommand{\TablaDeDatos}{\subsubsection{Tabla de datos}}

% Palabras clave para resaltar en el texto.
\newcommand{\PalabraClave}[1]{\medskip \textbf{#1}}

% Título para gráficos.
\newcommand{\TituloGrafico}[1]{\subsubsection{Gráfico: #1}}

% Negrita-Italica para resaltar variables.
\newcommand{\Variable}[1]{\textbf{\textit{#1}}}

% Subsubsection Interpretación.
\newcommand{\Interpretacion}{\subsubsection{Interpretación}}
